% ══════════════════════════════════════════════════════════════════════════════
% CHAPTER 2 — LITERATURE REVIEW  [revised]
% ══════════════════════════════════════════════════════════════════════════════

This chapter reviews the academic literature underpinning the
methodology and research questions of this thesis. It proceeds in
eight sections covering central bank communication theory
(Section~\ref{sec:lit_communication}), the text-as-data methodology
(Section~\ref{sec:lit_textasdata}), dictionary approaches to
monetary policy text (Section~\ref{sec:lit_dictionary}), machine
learning and transformer-based methods
(Section~\ref{sec:lit_ml}), the event-study identification strategy
(Section~\ref{sec:lit_eventstudy}), forward guidance
(Section~\ref{sec:lit_forward}), interest rate predictability
(Section~\ref{sec:lit_inertia}), and the gaps this thesis addresses
(Section~\ref{sec:lit_gaps}).
 


% ══════════════════════════════════════════════════════════════════════════════
\section{Central Bank Communication as a Policy Tool}
\label{sec:lit_communication}
% ══════════════════════════════════════════════════════════════════════════════
Modern monetary policy operates not only through policy-rate decisions,
but also through communication. The theoretical foundation is the
expectations channel: economic activity and asset prices depend not only
on the current short-term rate, but also on the expected future path of
policy and the central bank's reaction function \citep{woodford2005}.
If communication shapes these expectations, it becomes a policy
instrument in its own right.

This view is supported by \citet{blinder2008}, who survey the evidence
on central bank communication and argue that it can move financial
markets, improve the predictability of monetary policy decisions, and
contribute to the effectiveness of monetary policy. At the zero lower
bound, this role became particularly important, as forward guidance
allowed central banks to influence expectations about future policy when
further conventional rate cuts were constrained \citep{campbell2012}.

The institutional development of communication also differs across
central banks. The FOMC began issuing regular post-meeting statements in
February~1994, while the ECB has held press conferences since the start
of monetary union in 1999. These differences matter because the two
institutions communicate through different formats and event windows.
The empirical analysis therefore studies FOMC and ECB communication
separately, allowing the estimated communication--market relationship to
differ across institutional settings.

Central bank communication also provides guidance to financial markets.
When a central bank describes economic activity as ``strong'' or
inflation pressures as ``broad-based'', it is not only reporting current
conditions. It is also guiding markets about how the central bank is
likely to interpret those conditions and how policy may respond. This
signalling function is what makes text-based tone measurement useful:
specific words and phrases can reveal whether the central bank's
assessment points toward tightening, easing, or maintaining the current
policy stance. The challenge, addressed in
Section~\ref{sec:lit_dictionary}, is to identify which terms carry
genuine stance signals and which are neutral descriptions.


% ══════════════════════════════════════════════════════════════════════════════
\section{Text-as-Data in Economics and Finance}
\label{sec:lit_textasdata}
% ══════════════════════════════════════════════════════════════════════════════
The central premise of text-as-data research is that economically
relevant information is embedded in language and can be measured
systematically. Unlike standard numerical variables, however, text must
first be transformed into quantitative measures. This requires
researcher choices about document selection, cleaning, tokenisation,
feature construction, weighting, aggregation, and validation. These
choices matter because the resulting measure must correspond to the
economic concept it is intended to capture \citep{gentzkow2019}.

This thesis applies this framework to central bank communication. The
main object of interest is hawkish--dovish monetary-policy tone in FOMC
and ECB communication. The empirical question is whether this tone
explains financial market reactions beyond the monetary policy rate
surprise. The thesis therefore compares three types of tone measures.
The dictionary approach is transparent but relatively rigid: it counts
predefined hawkish and dovish terms, but may miss context such as the
difference between ``increasing asset purchases'' and ``decreasing asset
purchases''. The Random Forest approach is more flexible, as it can learn
nonlinear patterns in bigram TF-IDF features, but is less directly
interpretable. The RoBERTa-based approach, following \citet{shah2023},
uses contextual sentence-level classification into hawkish, dovish, and
neutral sentences, allowing the meaning of a phrase to depend on the
surrounding sentence. Comparing these methods evaluates the trade-off
between transparency, flexibility, and context sensitivity in measuring
central bank tone.

\subsection{Foundational Results}

General-purpose sentiment dictionaries are often poorly suited to
specialised economic and financial language. \citet{loughran2011}
show that many words classified as negative by the Harvard dictionary
are not negative in financial filings; terms such as ``liability'' and
``depreciation'' are often neutral accounting vocabulary rather than
negative sentiment. This illustrates a broader problem for
dictionary-based text measurement: the meaning of words depends on the
domain in which they are used.

The same concern applies to central bank communication. Monetary-policy
terms such as ``tightening'', ``slack'', ``accommodation'', or ``asset
purchases'' cannot be interpreted reliably with a generic sentiment
dictionary. Their meaning depends on the policy context and on whether
the sentence signals a more hawkish or dovish stance. This motivates the
use of a monetary-policy-specific dictionary in this thesis, as well as
data-driven benchmarks based on Random Forest and RoBERTa that can
detect relevant textual patterns without relying exclusively on manually
selected terms.


% ══════════════════════════════════════════════════════════════════════════════
\section{Dictionary-Based Approaches to Monetary Policy Text}
\label{sec:lit_dictionary}
% ══════════════════════════════════════════════════════════════════════════════

\subsection{Prior Dictionaries}
An early and influential approach to automated textual analysis of FOMC
statements is \citet{lucca2009}. Lucca and Trebbi develop a
computational measure of the hawkish or dovish content of FOMC
statements, with a particular focus on communication about future
interest rate decisions. Their results show that FOMC statements contain
information about future policy decisions and that this information is
reflected especially in longer-maturity Treasury yields. Their approach
therefore provides an important antecedent for this thesis, which also
asks whether textual communication contains information beyond the
current policy-rate decision.

\citet{apel2014} apply a dictionary-based approach to Riksbank minutes
and find that the textual content of the minutes helps predict future
policy-rate decisions. Their dictionary is based on adjective--noun
pairs, such as ``high inflation'' or ``weak growth'', which allows the
same economic concept to carry different policy implications depending
on the surrounding qualifier. This motivates the use of phrase-based
rather than purely word-based measurement in the present thesis.

\citet{picault2017} construct an ECB-specific weighted lexicon from
labelled sentences. Sentences are first classified as hawkish, neutral,
or dovish for monetary policy, and as positive, neutral, or negative for
the economic outlook. The authors then assign weights to words and
n-grams according to how frequently they appear in each class. For
example, an expression that occurs mainly in hawkish sentences receives a
hawkish weight, while an expression that occurs mainly in dovish
sentences receives a dovish weight. This weighted lexicon improves on
generic dictionaries by using vocabulary specific to ECB communication.
The dictionary developed in this thesis takes a different route: rather
than estimating weights from labelled sentences, it organises
hand-curated terms into four pre-specified semantic categories: policy
action, inflation assessment, employment and growth, and residual risk
language (\Cref{sec:lit_multidim}).

\citet{schmeling2019} study ECB press conferences and show that changes
in central bank tone have statistically and economically significant
effects on asset prices. A more positive ECB tone is associated with
higher stock prices and interest rates, and with lower credit spreads,
volatility risk premia, and implied volatility. Their results show that
communication tone can matter for financial markets even after
controlling for policy actions and conventional monetary-policy shocks.
This thesis builds on this event-based logic, but focuses specifically on
hawkish--dovish monetary-policy tone and compares dictionary, Random
Forest, and RoBERTa-based measures on a common FOMC and ECB corpus.

\Cref{tab:dict_comparison} summarises the major monetary-policy
text-scoring approaches alongside the present thesis.

% --------------------------------------------------------------------------
% Table 2.1 — Enhanced dictionary comparison
% Adapted from the Apel & Grimaldi (2014) Table 1 concept of showing
% dictionary construction dimensions side by side, extended to include
% all four benchmark studies plus this thesis.
% --------------------------------------------------------------------------
\begin{table}[htbp]
\centering
\caption{Comparison of monetary-policy text-scoring approaches}
\label{tab:dict_comparison}
\begin{threeparttable}
\scriptsize
\begin{tabularx}{\textwidth}{@{}
>{\raggedright\arraybackslash}p{2.2cm}
>{\raggedright\arraybackslash}p{1.3cm}
>{\raggedright\arraybackslash}p{1.5cm}
>{\raggedright\arraybackslash}X
>{\raggedright\arraybackslash}p{1.2cm}
>{\raggedright\arraybackslash}p{2.1cm}
@{}}
\toprule
Paper & Bank(s) & Period & Approach & N-grams? & Decomposition \\
\midrule

\citet{lucca2009}
  & FOMC
  & 1999--2008\tnote{a}
  & Search-engine semantic orientation
  & No
  & Single axis \\

\addlinespace
\citet{apel2014}
  & Riksbank
  & 2004--2011\tnote{b}
  & Adjective--noun phrase counts
  & Pairs\tnote{c}
  & Single axis \\

\addlinespace
\citet{hansen2016}
  & FOMC
  & 1998--2014
  & LDA topics and FAVAR
  & No
  & Two dimensions\tnote{d} \\

\addlinespace
\citet{hansen2018}
  & FOMC
  & 1976--2007
  & LDA on transcripts
  & No
  & Topic-level\tnote{e} \\

\addlinespace
\citet{picault2017}
  & ECB
  & 2006--2014\tnote{b}
  & Supervised n-gram weights
  & Yes
  & Two categories\tnote{f} \\

\addlinespace
\citet{schmeling2019}
  & ECB
  & 1999--2021
  & LM dictionary on press conferences
  & No
  & Single axis \\

\addlinespace
\citet{gorodnichenko2023}
  & FOMC
  & 2011--2019
  & BERT and voice tone
  & N/A\tnote{g}
  & Text versus voice \\

\midrule
\textbf{This thesis}
  & FOMC and ECB
  & 1994--2026
  & Hand-curated dictionary, Random Forest, and RoBERTa
  & Yes
  & Four categories\tnote{h} \\

\bottomrule
\end{tabularx}

\begin{tablenotes}\scriptsize
  \item[a] Exact start date to be verified against the paper.
  \item[b] Exact sample period to be verified against the paper's data section.
  \item[c] Adjective--noun pairs, such as ``high inflation'', are functionally similar to bigrams.
  \item[d] Economic conditions and forward guidance are identified through LDA topics, with effects traced through a FAVAR.
  \item[e] Latent topics emerge from the data; the paper focuses on how transparency affected FOMC deliberation.
  \item[f] Monetary-policy inclination and economic-outlook assessment, based on probabilistic n-gram weights learned from labelled sentences.
  \item[g] Contextual embeddings rather than dictionary-based term matching.
  \item[h] Policy action, inflation assessment, employment and growth, and residual risk language.
\end{tablenotes}

\end{threeparttable}
\end{table}

\subsection{Methodological Challenges in Dictionary Design}

Despite their transparency, dictionary approaches face several
limitations.

First, \emph{term selection} involves researcher judgment. Different
researchers may construct different hawkish and dovish word lists from
the same source material. This thesis addresses this limitation by using
explicit inclusion rules for the hand-curated dictionary and by
comparing the resulting scores with data-driven alternatives based on
Random Forest and RoBERTa.

Second, dictionaries are only partially sensitive to context. For
example, ``inflation is rising'' and ``inflation is not rising'' may
receive similar scores if the dictionary does not account for negation.
Matching phrases at the bigram and trigram level partly reduces this
problem by capturing expressions such as ``rising inflation'' or
``inflation pressures'', but fixed term lists cannot fully account for
negation, sentence structure, or changes in meaning across contexts.
This motivates the comparison with RoBERTa, which classifies sentences
using contextual language representations.

Third, dictionaries are static. The same word list is applied across a
period in which central bank communication changes substantially. Terms
such as ``data dependent'', ``forward guidance'', or ``average inflation
targeting'' illustrate how new policy vocabulary can emerge over time.
For this reason, dictionary-based scores should be interpreted as
transparent but imperfect approximations of monetary-policy tone, rather
than as complete measures of communication.

\subsection{Validation of Tone Measures}
 
A critical question for any tone measure is how to assess its validity.
The literature uses several approaches. \citet{picault2017} validate
their ECB dictionary by testing whether their indicators help explain
future ECB rate decisions in an augmented Taylor rule and whether they
are associated with stock market volatility on press-conference days.
They also compare their field-specific indicators with general-purpose
alternatives based on \citet{loughran2011} and \citet{apel2014}, showing
that the ECB-specific lexicon performs better.

\citet{lucca2009} use a different validation strategy by examining
whether their FOMC communication measure is reflected in financial
market reactions around monetary policy announcements. The present
thesis follows both validation logics. First, it examines whether
hawkish and dovish tone scores align with observed policy-rate
decisions. Second, it estimates market-reaction regressions to test
whether tone explains high-frequency financial market movements after
controlling for the monetary policy rate surprise. This combination
allows the analysis to assess both the policy relevance and the market
relevance of the proposed tone measures.
 
\subsection{Cross-Bank Dictionaries}
 
Most existing monetary-policy dictionaries are developed for a specific
institution, such as the FOMC in \citet{lucca2009} or the ECB in
\citet{picault2017}. Applying one dictionary across central banks is
therefore not straightforward, since institutions differ in mandates,
communication formats, and preferred terminology. For example, FOMC
communication places greater emphasis on labour-market conditions,
reflecting its dual mandate, while ECB communication is more strongly
centred on price stability and euro-area inflation dynamics.

The dictionary used in this thesis is designed to be applicable to both
institutions. It includes common monetary-policy terminology as well as
institution-specific expressions, such as ``second-round effects'' and
``below our aim'' for the ECB, and ``labor market slack'' or
``employment risks'' for the FOMC. The analysis therefore examines
whether a common cross-bank dictionary can capture hawkish--dovish tone
in both institutional settings.
 
\subsection{Multi-Dimensional Tone Measurement}
\label{sec:lit_multidim}
 
A common feature of monetary-policy dictionaries is that they often
collapse communication into a single hawkish--dovish score. While this
is useful for summarising overall tone, it can obscure which part of the
communication drives the signal. Central bank statements typically
combine several types of information: policy-action language, inflation
assessment, labour-market and real-activity commentary, and broader risk
or residual communication. A single index may therefore hide whether a
market reaction is associated with inflation language, policy guidance,
or real-side economic assessment.

Several studies suggest that this distinction matters. For example,
\citet{hansen2016} distinguish between communication about current
economic conditions and communication about the future path of policy,
while \citet{picault2017} separate monetary-policy inclination from
economic-outlook assessment in ECB press conferences. These approaches
show that different dimensions of central bank communication can carry
different information.

This thesis follows this logic by organising the dictionary into four
pre-specified semantic categories: policy action, inflation assessment,
employment and real-activity conditions, and a residual category for
other relevant monetary-policy language. The categories are not meant to
represent separate structural shocks. In practice, central bank language
often links several dimensions: wage pressures can signal both
labour-market strength and inflation risk, while weak demand can imply
both real-side deterioration and future policy accommodation. The purpose
of the decomposition is therefore to make the tone measure more
interpretable and to examine which broad dimension of communication is
most closely associated with each market response.
 
 % ══════════════════════════════════════════════════════════════════════════════
\section{Machine Learning and NLP in Central Bank Analysis}
\label{sec:lit_ml}
% ══════════════════════════════════════════════════════════════════════════════



\subsection{Supervised Classification}

Random Forests \citep{breiman2001} occupy a methodological middle
ground between rigid dictionary methods and more complex deep-learning
models. A Random Forest is an ensemble of decision trees, where each
tree is trained on a bootstrap sample of the data and considers only a
random subset of features at each split. Averaging across many trees
reduces the instability and overfitting risk of a single decision tree,
while still allowing some interpretation through feature-importance
measures.

In this thesis, the Random Forest is used as a supervised benchmark for
dictionary-based tone measurement. Instead of relying on pre-specified
hawkish and dovish terms, the model learns which bigram features are
most informative for classifying central bank communication. Out-of-bag
(OOB) predictions provide an internal validation measure: each
observation is predicted only by trees that were not trained on that
observation, yielding a cross-validation-like estimate without requiring
a separate holdout sample \citep{breiman2001}.

\subsection{Transformer-Based Approaches}
\label{sec:lit_transformers}

The introduction of transformer architectures \citep{vaswani2017} and
pre-trained language models such as BERT \citep{devlin2019} has changed
the frontier of text classification. Unlike bag-of-words methods,
transformers process words in context. This is important for central
bank communication because the same term can carry different meanings
depending on the surrounding sentence. For example, ``raised'' in
``the Committee raised rates'' conveys a different signal than in
``concerns were raised about downside risks''. Contextual models are
therefore better suited than fixed dictionaries to handling ambiguity,
negation, and sentence-level meaning.

\citet{shah2023} introduce the \emph{Trillion Dollar Words} dataset and
formulate hawkish--dovish--neutral classification as a sentence-level
task for FOMC communication. They compare transformer-based models with
simpler baselines and show that fine-tuned RoBERTa specifications
perform strongly in classifying monetary-policy stance. Their
contribution is important because it shifts hawkish--dovish measurement
from document-level word counts toward contextual sentence-level
classification.

This thesis uses a fine-tuned RoBERTa-based sentence classifier as a
third benchmark against the dictionary and Random Forest measures. The
purpose is not to treat RoBERTa as the preferred measure, but to test
whether a contextual sentence-level model captures information that is
missed by transparent dictionary counts and document-level
machine-learning features. Since transformer predictions depend on the
domain and quality of the fine-tuning data, the resulting scores are
interpreted as a benchmark rather than as a definitive measure of
central bank tone. The sentence-level aggregation procedure is described
in Chapter~\ref{ch:methodology}.

\subsection{Classification Challenges and Evaluation}

Supervised classification of central bank communication faces two main
methodological challenges in this thesis. First, policy-rate decisions
are imbalanced: unchanged policy rates occur more frequently than hikes
or cuts. A classifier can therefore achieve apparently high accuracy by
predicting the majority class too often, while failing to identify
tightening or easing episodes. For this reason, the evaluation does not
rely only on overall accuracy, but also considers metrics that are more
informative under class imbalance, such as balanced accuracy and
class-specific performance.

Second, the textual feature space is high-dimensional. Even when the
analysis is restricted to bigrams, the number of potential TF-IDF
features can become large relative to the number of central bank
documents. This creates both an overfitting problem and a computational
problem, since estimating Random Forest models on a very large sparse
feature space is slow and unstable. The thesis therefore uses
Lasso-based feature pre-selection before estimating the Random Forest.
This reduces the dimensionality of the bigram space, improves
computational efficiency, and limits the risk that the model fits noise
rather than informative textual patterns.

The methodological details of bigram construction, TF-IDF weighting,
Lasso-based pre-screening, Random Forest estimation
are described in Chapter~\ref{ch:methodology}.

\subsection{Model Comparison: Dictionary vs.\ Machine Learning
vs.\ Transformer}

A distinctive feature of this thesis is the direct comparison of three
tone-measurement approaches on the same corpus: a dictionary-based
measure, a Random Forest measure, and a RoBERTa-based sentence-level
measure. This comparison is useful because the methods differ in what
they can capture. The dictionary is transparent but relies on fixed
terms. The Random Forest can learn nonlinear patterns in bigram TF-IDF
features, but is less directly interpretable. RoBERTa can account for
sentence-level context, but its predictions depend on the quality and
domain of the fine-tuning data.

The empirical comparison is implemented through horse-race regressions
in which all three tone scores enter the same market-reaction
regression. If one measure fully captures the relevant information in
the others, only that measure should remain significant. If several
measures retain explanatory power, this suggests that they capture
different aspects of central bank tone. The purpose is therefore not only
to rank the methods, but also to assess whether transparent,
data-driven, and contextual approaches contain complementary
market-relevant information.

In addition, the thesis reports robustness specifications in which the
tone coefficients are estimated under alternative controls, including
chair and era fixed effects. This fixed-effects ladder examines whether
the estimated communication effects are robust to differences across
central bank leadership periods and broader monetary policy regimes.


% ══════════════════════════════════════════════════════════════════════════════
\section{Event Studies and the Identification of Market Reactions}
\label{sec:lit_eventstudy}
% ══════════════════════════════════════════════════════════════════════════════

\subsection{The Target and Path Components of Monetary Policy Announcements}

\citet{gurkaynak2005} show that financial market reactions to FOMC
announcements cannot be explained by the surprise change in the current
policy rate alone. Instead, they identify two components of monetary
policy announcements: a target factor, which captures news about the
current policy rate, and a path factor, which captures news about the
expected future path of policy. The path factor is closely associated
with FOMC statements and has particularly strong effects on longer-term
Treasury yields.

A central implication of the monetary-surprise literature is that
hawkishness is not equivalent to a realised rate increase. Financial
markets respond to the unexpected component of policy. A 25 basis point
rate increase may therefore be market-neutral if it was fully priced in
before the announcement, hawkish if markets expected a smaller increase,
or dovish if markets expected a larger increase. This distinction is
particularly important for text-based measurement, because statements
following rate hikes mechanically contain restrictive language such as
``rate increase'' or ``tightening'', even when the policy action itself
was anticipated. The literature therefore motivates separating realised
policy-decision stance from market-surprise stance.

This distinction motivates the regression design in this thesis. The
monetary policy rate surprise, measured by MP1 for the FOMC and the
corresponding OIS surprise for the ECB, is included as a control for the
unexpected current-rate component of the announcement. The coefficient
on \netscore{} is therefore interpreted as the association between
central bank tone and market reactions after controlling for the
mechanical surprise in the policy-rate decision.

\subsection{The Information Effect}

Central bank announcements may affect markets not only because they
change the expected path of policy, but also because they reveal the
central bank's assessment of the economy. This is known as the central
bank information effect. The key idea is that a policy announcement can
contain two types of news: news about monetary policy itself and news
about economic fundamentals.

\citet{nakamura2018} show that expected output growth rises after
unexpected monetary tightenings around FOMC announcements. This is
difficult to reconcile with a pure tightening shock, which would normally
be expected to lower future output. One interpretation is that a rate
increase can also signal that the Federal Reserve views the economy as
stronger than markets previously believed.

\citet{jarocinski2020} formalise this distinction using the
high-frequency co-movement of interest rates and stock prices around
central bank announcements. A pure monetary policy tightening raises
interest rates and lowers stock prices, reflecting the standard
discount-rate channel. By contrast, a positive central bank information
shock raises both interest rates and stock prices, because the
announcement reveals a stronger economic outlook.

This distinction is important for interpreting the equity regressions in
this thesis. A hawkish tone coefficient need not be negative for stock
prices. If hawkish communication mainly signals tighter future policy,
equities should fall. If it instead signals stronger economic
fundamentals, equities may rise. The expected sign of the tone
coefficient is therefore ambiguous for equities, while it is more
directly interpretable for interest-rate outcomes.

\subsection{Market-Specific Reactions}

The expected market response to monetary-policy communication differs
across asset classes. For equities, the standard discount-rate channel
implies that hawkish surprises should lower stock prices, while dovish
surprises should raise them. Consistent with this channel,
\citet{bernanke2005} find that an unanticipated 25~basis point cut in
the federal funds rate target is associated with an increase of roughly
1\% in broad stock indexes.

Exchange rates provide a second sign check. A more hawkish domestic
monetary-policy signal should, all else equal, appreciate the domestic
currency by raising expected interest-rate differentials. In this thesis,
EUR/USD is therefore interpreted differently across central banks: a
hawkish FOMC communication shock is expected to lower EUR/USD through
dollar appreciation, while a hawkish ECB communication shock is expected
to raise EUR/USD through euro appreciation.

Bond yields provide the most direct link to central bank communication.
\citet{gurkaynak2005} show that FOMC announcements contain both a
current-rate component and a future-path component, with the latter being
particularly important for longer-term Treasury yields. This motivates
examining yield reactions across the curve rather than focusing only on
the shortest maturity. In this thesis, the yield outcomes therefore test
whether hawkish--dovish tone is associated not only with immediate
policy-rate expectations, but also with broader revisions to the expected
future path of monetary policy.

\subsection{Monetary Policy Databases}

For the FOMC, intraday market reactions are taken from the U.S.\
Monetary Policy Event-Study Database of \citet{acosta2025}. The
database reports high-frequency changes in financial variables around
separate FOMC communication windows, including the statement window, the
press-conference window, and a broader monetary-event window. For the
ECB, the Euro Area Monetary Policy Database of \citet{altavilla2019}
provides analogous high-frequency asset-price changes around the press
release and press-conference windows.

This window separation is central to the empirical design. It allows the
analysis to match each textual corpus to the corresponding market
reaction window: FOMC statements to the FOMC statement window, FOMC press
conferences to the FOMC press-conference window, ECB policy statements
to the ECB press-release window, and ECB press conferences to the ECB
press-conference window. The resulting dependent variables are therefore
not daily asset-price changes, but high-frequency changes in yields,
equity prices, and exchange rates measured around specific monetary
policy communication events.



% ══════════════════════════════════════════════════════════════════════════════
\section{Forward Guidance and Forward-Looking Communication}
\label{sec:lit_forward}
% ══════════════════════════════════════════════════════════════════════════════

\



\subsection{Press Conferences as Information-Rich Channels}

Press conferences are a particularly rich source of central bank
communication because they contain both prepared remarks and less
scripted responses to journalists' questions. Compared with
post-meeting statements, this format can produce more varied language
and may reveal how the Chair or President interprets incoming data,
policy trade-offs, and risks to the outlook.

Recent evidence suggests that press conferences contain market-relevant
information beyond the written policy statement. \citet{gorodnichenko2023}
analyse the vocal tone of Federal Reserve Chairs during FOMC press
conferences and show that positive vocal tone is associated with
significant increases in share prices, even after controlling for Federal
Reserve actions and the sentiment of policy texts. Their results imply
that not only the content of communication, but also the way it is
delivered, can affect financial markets.

This thesis does not analyse vocal tone. Instead, it focuses on the
textual content of statements and press conferences. The evidence from
\citet{gorodnichenko2023} nevertheless motivates treating press
conferences as a distinct communication channel. Accordingly, the
empirical analysis studies press conferences separately from statements
and retains only the central banker's responses, excluding journalist
questions from the textual corpus.


% ══════════════════════════════════════════════════════════════════════════════
\section{Interest Rate Predictability and Policy Inertia}
\label{sec:lit_inertia}

The second research question asks whether central bank language contains
predictive information about near-term monetary policy outcomes. The
forecasting exercise focuses on two types of targets. First, future rate
changes measure the realised policy decision at the next meeting and two
meetings ahead. Second, future two-year yield changes measure how market
expectations about the near-term policy path evolve over the same
horizons.

This distinction is important. A central bank may use today’s statement
to prepare markets for a future rate change before the decision is
actually implemented. In that case, communication may affect two-year
yields immediately, even if the realised policy rate changes only at a
later meeting. Conversely, a model may help predict the next rate
decision without explaining yield movements if the decision was already
fully priced in by markets. The thesis therefore studies both realised
rate decisions and two-year yield changes.

The two-meeting horizon reflects the gradual nature of monetary policy.
Central banks often adjust rates in sequences and communicate the likely
direction of policy before taking the next step. A statement may
therefore contain information not only about the next meeting, but also
about whether the current policy path is likely to continue or reverse.
For this reason, RQ2 uses both one-meeting-ahead and two-meeting-ahead
forecasting targets for rate changes and two-year yields.

\subsection{Policy Inertia}

Rate decisions are persistent. Central banks often keep rates unchanged
for several meetings and, when they do adjust policy, rate changes tend
to occur in tightening or easing sequences. In the present sample, hold
decisions account for the majority of meetings in both the FOMC and ECB
corpora (\Cref{tab:rate_distribution}). This persistence creates an
important benchmark for the forecasting exercise: a model may appear to
predict future rate decisions simply because it captures the current
policy direction.

The RQ2 design therefore evaluates whether central bank language adds
predictive information beyond simple policy inertia. For rate-decision
targets, the current meeting's rate change is included as a control when
assessing the expanding-window forecasts. For yield targets, the current
two-year yield change provides the corresponding market-based inertia
benchmark. The key question is therefore whether textual measures improve
forecasts of future rate decisions and two-year yield changes beyond
what is already implied by the current policy move or market reaction.

\subsection{Communication as a Signaling Mechanism}

Why should communication at meeting~$t$ help forecast policy or yield
changes at meeting~$t+1$ or~$t+2$? One possibility is that the central
bank uses forward-looking language to signal the likely direction of
future policy, thereby preparing markets for upcoming actions. Another
possibility is that the language summarises the information set on which
future policy decisions are based. For example, if a statement describes
inflation as persistently above target and labour-market conditions as
tight, future tightening may become more likely because both the current
language and the later decision reflect the same underlying assessment.

This distinction is important for interpreting RQ2. Predictive power
does not imply that communication mechanically causes future policy
changes. Rather, it means that the language used at meeting~$t$ contains
information that helps forecast the realised rate decision or the change
in two-year yields at subsequent meetings. The two-year yield target is
included because it reflects market expectations about the near-term
policy path: yields may adjust even if the next realised rate decision is
unchanged.

\citet{apel2014} provide evidence for this policy-prediction channel by
showing that textual content in Riksbank minutes helps predict future
rate decisions. This thesis extends the logic to FOMC and ECB statements
and press conferences, and separates realised policy outcomes from
market-based expectations by forecasting both future rate changes and
future two-year yield changes.

% ══════════════════════════════════════════════════════════════════════════════
\section{Gaps in the Literature and Contribution}
\label{sec:lit_gaps}
% ══════════════════════════════════════════════════════════════════════════════

The literature reviewed above leaves five gaps that this thesis
addresses.

\textbf{Gap~1: Cross-bank, long-sample evidence.} Much of the
existing literature focuses on a single central bank, such as
\citet{lucca2009} for the FOMC, \citet{picault2017} for the ECB, or
\citet{schmeling2019} for ECB press conferences. This thesis applies a
consistent empirical pipeline to both the FOMC and the ECB over a long
sample period. The FOMC sample covers post-meeting statements from
1994 to 2026, while the ECB sample covers communication from 1999 to
2026. This allows the analysis to compare the communication--market
relationship across two major central banks and across different
monetary policy regimes.

\textbf{Gap~2: Multi-dimensional tone measurement.} Existing
dictionary approaches often summarise communication using a single
hawkish--dovish score. While useful, a single score can hide which part
of the communication drives the signal. This thesis therefore decomposes
tone into four broad semantic categories: policy action, inflation
assessment, employment and real-activity conditions, and a residual
category for other relevant monetary-policy language. The purpose is not
to construct separate structural shocks, but to examine which dimensions
of communication are most closely associated with different market
responses.

\textbf{Gap~3: Direct comparison of measurement methods.} Most studies
use either a dictionary-based approach or a machine-learning approach.
This thesis compares three methods on the same corpus: a hand-curated
dictionary, a Random Forest model, and a RoBERTa-based sentence-level
classifier. Horse-race regressions, in which all three scores enter the
same specification, test whether the methods capture overlapping or
independent market-relevant information.

\textbf{Gap~4: Forward-looking communication.} Forward guidance implies
that markets should react not only to descriptions of current economic
conditions, but also to language about the expected future path of
policy. This thesis therefore separates forward-looking from
current/backward-looking sentences using linguistic markers of futurity,
projection, commitment, and conditionality. The resulting dictionary and
RoBERTa-based forward-looking scores are used to examine whether
forward-looking communication carries distinct information for market
reactions and forecasting performance.

\textbf{Gap~5: Multi-horizon forecasting.} Existing work on textual
prediction of policy decisions often focuses on the next meeting. This
thesis uses an expanding-window design to forecast four near-term
targets: the rate decision at the next meeting, the rate decision two
meetings ahead, the two-year yield change at the next meeting, and the
two-year yield change two meetings ahead. This design tests whether
central bank language contains predictive information beyond simple
policy inertia and beyond the immediate next-meeting horizon.